\section{Latent Clustering of Surrogate Policies}
\label{sec:clustering_agent_representation_method}

Following the work of \cite{zahavy2017graying}, we analyze the \gls{ls} through clustering techniques, placing our approach within the broader field of Deep Clustering.
This design choice is motivated by our rejection of the \SuperpositionHypothesisName and, more generally, of any view that assumes a linearly organized \gls{cs}, as discussed in Section~\ref{sec:clustering}.

The \gls{car} methodology is composed of two main stages.
The first stage trains multiple independent \gls{nn}, called surrogate policies, which implement the same decision policy.
The second stage clusters their \gls{lr} in order to identify the underlying conceptual organization.
Once these two stages have been completed, the resulting clustering can be analyzed from two complementary perspectives.
First, the consistency of the resulting clusters across independent training runs can be assessed to determine whether the identified clusters are stable.
Second, the semantic content of the clusters can be examined to determine whether they correspond to meaningful concepts.
Figure~\ref{fig:illustration_method} provides a visual overview of the two stages, from the imitation of the initial policy to the clustering of the surrogate representations.

\begin{figure}[t]
    \centering
    \resizebox{\textwidth}{!}{
    \begin{tikzpicture}
        \node (observation) [minimum height=1.5cm, minimum width=2.5cm, align=center] at (0, 0) {\hl{Observation} \\ \large $O$};

        \node[draw, rectangle, text centered, minimum height=1cm, minimum width=2.5cm, align=center, rounded corners=0.2cm] (policy) at ($(observation) + (3cm, 1.35cm)$) {\scalebox{2}{$\textcolor{mygreen}{\pi}$}};

        \node[align=center] (policy_text) at ($(policy) + (0, 1cm)$) {Initial Policy};

        \node[draw, rectangle, text centered, minimum height=1cm, minimum width=2.5cm, align=center, rounded corners=0.2cm] (surrogate_policy) at ($(observation) + (3cm, -1.35cm)$) {\scalebox{2}{$\textcolor{myred}{\pi}^{\textcolor{myred}{'}}_{\textcolor{myred}{\Theta}}$}};

        \node[align=center] (surrogate_policy_text) at ($(surrogate_policy) + (0, 1cm)$) {Surrogate Policy};

        \draw[->] (observation) |- (policy);
        \draw[->] (observation) |- (surrogate_policy);

        \def\q{5};
        \def\B{4};
        \def\S{5};
        \def\Bs{1.0};
        \def\Ss{1.25};
        \def\xmax{\S+3.2*\Ss};
        \def\ymin{{-0.15*gauss(\B,\B,\Bs)}};

        \begin{axis}[
            name=policy_distribution,
            at={($(policy) + (5cm, 0cm)$)},
            anchor=center,
            every axis plot post/.append style={
                mark=none,domain={-0.05*(\xmax)}:{1.08*\xmax},samples=\numberOfSamplesToPlotCurve,smooth
            },
            xmin=0,
            xmax=\xmax,
            ymin=0,
            ymax={1.1*gauss(\B,\B,\Bs)},
            axis lines=middle,
            axis line style=thick,
            enlargelimits=upper, % extend the axes a bit to the right and top
            ticks=none,
            x=0.5cm,
            y=5cm,
            scale=0.7,
            xlabel={$\mathcal{A}$},
            xlabel style={at={(axis description cs:0.85,0.2)}, anchor=west},
            % ylabel={probability},
            ylabel style={at={(axis description cs:0,1.02)},anchor=south},
            title={action distribution},
            title style={at={(axis description cs:0.5,-0.15)},anchor=north},
        ]

        % PLOTS
        \addplot[name path=A,thick, mygreen] {gauss(x,\B,\Bs)} node[pos=0.5, above, xshift=-0.8cm, yshift=-0.2cm] {$\pi\textcolor{black}{(O)}$};
        \end{axis}

        \begin{axis}[
            name=surrogate_policy_distribution,
            at={($(surrogate_policy) + (5cm, 0cm)$)},
            anchor=center,
            every axis plot post/.append style={
                mark=none,domain={-0.05*(\xmax)}:{1.08*\xmax},samples=\numberOfSamplesToPlotCurve,smooth
            },
            xmin=0,
            xmax=\xmax,
            ymin=0,
            ymax={1.1*gauss(\B,\B,\Bs)},
            axis lines=middle,
            axis line style=thick,
            enlargelimits=upper, % extend the axes a bit to the right and top
            ticks=none,
            x=0.5cm,
            y=5cm,
            scale=0.7,
            xlabel={$\mathcal{A}$},
            xlabel style={at={(axis description cs:0.85,0.2)}, anchor=west},
            % ylabel={probability},
            ylabel style={at={(axis description cs:0,1.02)},anchor=south},
            title={action distribution},
            title style={at={(axis description cs:0.5,-0.15)},anchor=north},
        ]

        % PLOTS
        \addplot[name path=B,thick,myred] {gauss(x,\S,\Ss)} node[pos=0.5, above, xshift=1.2cm, yshift=-0.3cm] {$\pi^{'}_{\Theta}\textcolor{black}{(O)}$};
        \end{axis}

        \draw[->] (policy) -- ($(policy_distribution.west) + (-0.2cm, 0)$);
        \draw[->] (surrogate_policy) -- ($(surrogate_policy_distribution.west) + (-0.2cm, 0)$);

        \begin{axis}[
            name=distribution_match,
            at={(13.5cm, 0cm)},
            anchor=center,
            every axis plot post/.append style={
                mark=none,domain={-0.05*(\xmax)}:{1.08*\xmax},samples=\numberOfSamplesToPlotCurve,smooth
            },
            xmin=0,
            xmax=\xmax,
            ymin=0,
            ymax={1.1*gauss(\B,\B,\Bs)},
            axis lines=middle,
            axis line style=thick,
            enlargelimits=upper, % extend the axes a bit to the right and top
            ticks=none,
            x=0.5cm,
            y=5cm,
            scale=0.8,
            xlabel={$\mathcal{A}$},
            xlabel style={at={(axis description cs:0.85,0.2)}, anchor=west},
            % ylabel={probability},
            ylabel style={at={(axis description cs:0,1.02)},anchor=south},
            title={\shortstack{difference between \\[0.08cm] distributions}},
            title style={at={(axis description cs:0.5,-0.15)},anchor=north},
        ]

        % PLOTS
        \addplot[name path=A,thick,mygreen] {gauss(x,\B,\Bs)} node[pos=0.5, above, xshift=-0.8cm, yshift=-0.2cm] {$\pi\textcolor{black}{(O)}$};
        \addplot[name path=B,thick,myred] {gauss(x,\S,\Ss)} node[pos=0.5, above, xshift=1.2cm, yshift=-0.3cm] {$\pi^{'}_{\Theta}\textcolor{black}{(O)}$};

        \addplot[opacity=0.2] fill between[of=A and B, soft clip={domain=0:15}];
        \end{axis}

        \node[align=center] at ($(distribution_match.north) + (0, 1.25cm)$) {$\textcolor{myred}{\pi}^{\textcolor{myred}{'}}_{\textcolor{myred}{\Theta}}$ is trained to be as close as \\ possible of $\textcolor{mygreen}{\pi}$ for all $O$ in  order \\ to reproduce its behavior.};

        \draw[->] ($(policy_distribution.east) + (0.2cm, 0)$) -- ($(policy_distribution.east) + (0.5cm, 0)$) |- ($(distribution_match.west) + (-0.2cm, 0)$);
        \draw[->] ($(surrogate_policy_distribution.east) + (0.2cm, 0)$) -- ($(surrogate_policy_distribution.east) + (0.5cm, 0)$) |- ($(distribution_match.west) + (-0.2cm, 0)$);

        \begin{scope}[xshift=8cm, yshift=-5cm, scale=0.8, transform shape]
            \node[draw, minimum width=4cm, minimum height=4cm] (box_2) at (0, 0) {};

            \begin{scope}
                \clip (box_2.south west) rectangle (box_2.north east);

                \node[draw, dashed, minimum width=1.5cm, circle, minimum height=1.5cm] (circle_1) at (-1.25cm, -1.25cm) {};

                \node[draw, dashed, minimum width=1.5cm, circle, minimum height=1.5cm] (circle_2) at (0cm, 1.75cm) {};

                \node[draw, dashed, minimum width=1.5cm, circle, minimum height=1.5cm] (circle_3) at (1.75cm, 0cm) {};
            \end{scope}

            \foreach \x/\y in {
                -1.5/-0.75,
                -1.25/-1.5,
                -1/-1.75,
                -1.75/-1.25,
                -0.75/-1.25,
                -1/-1,
                0.5/1.5,
                -0.25/1.25,
                0.0/1.50,
                -0.5/1.5,
                0/1.25,
                1.5/-0.5,
                1.5/0.5,
                1.25/0.25,
                1.5/0,
                1.25/-0.25
            }
            {
                \node[] (p1) at (\x cm,\y cm) {\textcolor{myblue}{$\times$}};
            }
        \end{scope}
        \draw[->, dashed, preaction={draw=white,line width=6pt}] (surrogate_policy) |- ($(box_2.west) + (-0.2cm, 0)$);

        \node[draw, rectangle, text centered, align=center, dashed, fill=white, inner sep=1.5mm] at ($(surrogate_policy) - (0, 2.0cm)$) {Projection obtained through \\ the function $\textcolor{myblue}{f_{enc}}$, which is \\ a subcomponent of $\textcolor{myred}{\pi}^{\textcolor{myred}{'}}_{\textcolor{myred}{\Theta}}$.};

        \node[align=center] at ($(box_2.east) + (2.75cm, 0)$) {Once $\textcolor{myred}{\pi}^{\textcolor{myred}{'}}_{\textcolor{myred}{\Theta}}$ is trained,\\ observations are projected into $\textcolor{myblue}{\mathcal{Z}}$,\\ clustered automatically,\\ and interpreted for semantics.};

        \node[align=center] at ($(box_2.south) + (0, -0.3cm)$) {latent space $\textcolor{myblue}{\mathcal{Z}}$ of $\textcolor{myred}{\pi}^{\textcolor{myred}{'}}_{\textcolor{myred}{\Theta}}$};
    \end{tikzpicture}
    }
    \caption{Illustration of the \methodName method.}
    \label{fig:illustration_method}
\end{figure}


This section is organized into three parts.
The first explains why multiple surrogate policies are used to study the stability of the representations.
The second describes how these policies are trained to match the initial policy as closely as possible.
The third presents how their \gls{lr} are clustered to identify the \gls{cs}.

\subsection{Multiple surrogate policies}

The \gls{car} method was developed to test the \RepresentationConvergenceHypothesisName.
This requires training several \gls{nn} independently and comparing their \gls{lr}.
Consistent similarities in the resulting \gls{cs} would provide empirical support for this hypothesis.

In supervised learning, this protocol is straightforward since several models can be independently trained on the same target function.
\gls{rl} introduces an additional difficulty.
Because of stochastic exploration, independently trained agents are not guaranteed to converge toward the same policy.
Even if they achieve comparable performance, they may solve the task using different strategies.
This makes it difficult to determine whether differences in the \gls{ls} result from different representations of the same strategy or from differences in the learned strategies themselves.

To overcome this limitation, we consider an already trained initial policy.
This policy is then used as a target for imitation learning \cite{pmlr-v15-ross11a}.
In this process, several surrogate \gls{nn} are independently trained in a supervised manner to reproduce its behavior.
This makes it possible to investigate whether equivalent policies also exhibit equivalent conceptual organizations.

It is important to emphasize that the use of surrogate policies does not make \gls{car} a surrogate explainability method, as defined in Subsection~\ref{subsec:explainable_artificial_intelligence}.
The surrogate policies remain deep \gls{nn}.
Their role is not to simplify the decision process but to provide multiple independent realizations of the same policy function, enabling the study of representation convergence.

\subsection{Training surrogate policies}

Throughout the remainder of this manuscript, the initial policy we aim to explain is denoted as $\InitialPolicy$.
This initial policy is obtained without imposing any constraints on its origin.
It may result from classical or deep reinforcement learning \cite{sutton2018reinforcement}, alpha-beta algorithm \cite{knuth1975analysis}, or even from demonstrations provided by humans \cite{pmlr-v9-ross10a}.
The crucial requirement is the availability of a dataset containing a sufficient number of projections from the observation space $\mathcal{O}$ to the action space $\mathcal{A}$ induced by $\InitialPolicy$ to train a \gls{nn}.

The replicated policies are referred to as surrogate policies.
To assess the stability of the learned representations, $n$ surrogate policies are trained independently, denoted by $\SurrogatePolicy{1}, \dots, \SurrogatePolicy{n}$.
All surrogate policies share the same architecture.
It consists of a \gls{nn} with two components, an encoder $\textcolor{myblue}{f_{enc}}$ and a decoder $f_{dec}$, highlighted by the blue and black dashed rectangles, respectively, in Figure~\ref{fig:surrogate_policy_architecture}.
The encoder maps the observation $O$ to the embedding, and the decoder reconstructs the parameters of the action distribution from this embedding.
For a given surrogate policy $\SurrogatePolicy{i}$, parameterized by $\Theta_i$, its encoder and decoder are written $\textcolor{myblue}{f_{enc_{i}}}$ and $f_{dec_{i}}$.

\begin{figure}[t]
    \centering
    \resizebox{\textwidth}{!}{
    \begin{tikzpicture}[x=3cm,y=1cm]
    \large
    \message{^^JDeep convolution neural network}
    \readlist\Nnod{4,5,6,5,3}
    \def\nstyle{1}
    \tikzset{
    node 1/.style={thick,circle,draw,minimum size=18,inner sep=0.5,outer sep=0.6},
    }

    \message{^^J  Layer}
    \foreachitem \N \in \Nnod{
        \def\lay{\Ncnt}
        \pgfmathsetmacro\prev{int(\Ncnt-1)}
        \message{\lay,}
        \foreach \i [evaluate={\y=\N/2-\i+0.5; \x=\lay; \n=\nstyle;}] in {1,...,\N}{ % loop over nodes
            \ifboolexpr{test {\ifnumcomp{\lay}{=}{3}}}
            {
                \node[node \n,outer sep=0.6] (N\lay-\i) at (\x,\y) {\textcolor{myblue}{$e_{\lay,\i}$}};
            }
            {
                \node[node \n,outer sep=0.6] (N\lay-\i) at (\x,\y) {};
            }

            \ifnum\lay>1
            \foreach \j in {1,...,\Nnod[\prev]}{ % loop over nodes in previous layer
                \draw[] (N\prev-\j) -- (N\lay-\i);
            }
            \fi
        }
    }

    \node[draw, rectangle, fit=(N1-1)(N3-6)(N3-1)(N5-1), inner sep=0.25cm, dashed, rounded corners=0.3cm, color=mygreen] (policy) {};
      \node[] (policy_label) at ($(policy.north) + (0, 0.5cm)$) {\scalebox{2.5}{$\textcolor{mygreen}{\pi}$}};

    % Observation
    \node[minimum height=1cm, minimum width=1cm] (observation_1) at ($(N1-1) + (-1.75cm, 0)$) {\scalebox{1.25}{$o_1$}};
    \node[minimum height=1cm, minimum width=1cm] (observation_2) at ($(N1-2) + (-1.75cm, 0)$) {\scalebox{1.25}{$o_2$}};
    \node[minimum height=1cm, minimum width=1cm] (observation_3) at ($(N1-3) + (-1.75cm, 0)$) {\scalebox{1.25}{$o_3$}};
    \node[minimum height=1cm, minimum width=1cm] (observation_4) at ($(N1-4) + (-1.75cm, 0)$) {\scalebox{1.25}{$o_4$}};

    \node[minimum height=1cm, minimum width=1cm, align=center] (observation_label) at ($(observation_1) + (0, 1.25cm)$) {\hl{Observation}};

    \draw[thick] (observation_1.north east) -- ($(observation_1.north east) + (-0.1cm, 0)$);
    \draw[thick] (observation_1.north east) -- (observation_4.south east);
    \draw[thick] (observation_4.south east) -- ($(observation_4.south east) + (-0.1cm, 0)$);

    \draw[thick] (observation_1.north west) -- ($(observation_1.north west) + (0.1cm, 0)$);
    \draw[thick] (observation_1.north west) -- (observation_4.south west);
    \draw[thick] (observation_4.south west) -- ($(observation_4.south west) + (0.1cm, 0)$);

    \draw[->] ($(observation_1.east) + (-0.15cm, 0)$) -- ($(N1-1.west) + (-0.0cm, 0)$);
    \draw[->] ($(observation_2.east) + (-0.15cm, 0)$) -- ($(N1-2.west) + (-0.0cm, 0)$);
    \draw[->] ($(observation_3.east) + (-0.15cm, 0)$) -- ($(N1-3.west) + (-0.0cm, 0)$);
    \draw[->] ($(observation_4.east) + (-0.15cm, 0)$) -- ($(N1-4.west) + (-0.0cm, 0)$);

    % Distribution
    \node[minimum height=1cm, minimum width=1cm] (distribution_1) at ($(N5-1) + (1.75cm, 0)$) {\scalebox{1.25}{$d_1$}};
    \node[minimum height=1cm, minimum width=1cm] (distribution_2) at ($(N5-2) + (1.75cm, 0)$) {\scalebox{1.25}{$d_2$}};
    \node[minimum height=1cm, minimum width=1cm] (distribution_3) at ($(N5-3) + (1.75cm, 0)$) {\scalebox{1.25}{$d_3$}};

    \draw[thick] (distribution_1.north east) -- ($(distribution_1.north east) + (-0.1cm, 0)$);
    \draw[thick] (distribution_1.north east) -- (distribution_3.south east);
    \draw[thick] (distribution_3.south east) -- ($(distribution_3.south east) + (-0.1cm, 0)$);

    \draw[thick] (distribution_1.north west) -- ($(distribution_1.north west) + (0.1cm, 0)$);
    \draw[thick] (distribution_1.north west) -- (distribution_3.south west);
    \draw[thick] (distribution_3.south west) -- ($(distribution_3.south west) + (0.1cm, 0)$);

    \draw[->] ($(N5-1.east) + (0.0cm, 0)$) -- ($(distribution_1.west) + (0.2cm, 0)$);
    \draw[->] ($(N5-2.east) + (0.0cm, 0)$) -- ($(distribution_2.west) + (0.2cm, 0)$);
    \draw[->] ($(N5-3.east) + (0.0cm, 0)$) -- ($(distribution_3.west) + (0.2cm, 0)$);

    \node[minimum height=1cm, minimum width=1cm, align=center] (distribution_label) at ($(distribution_1) + (0, 1.5cm)$) {Parameters \\ for the \\ \hl{Distribution}};

    % Embedding
    \node[align=center] (embedding_label) at ($(N3-6) + (0cm, -2cm)$) {\hl{Embedding}};

    \draw[thick, dashed, ->] ($(N3-6) + (0cm, -0.5cm)$) -- ($(embedding_label) + (0cm, 0.5cm)$);

    \node[] (embedding) at ($(embedding_label) + (0, -2cm)$) {
        $\begin{bmatrix}
            \textcolor{myblue}{e_{3,1}} \\
            \textcolor{myblue}{e_{3,2}} \\
            \textcolor{myblue}{e_{3,3}} \\
            \textcolor{myblue}{e_{3,4}} \\
            \textcolor{myblue}{e_{3,5}} \\
            \textcolor{myblue}{e_{3,6}}
        \end{bmatrix}$
    };

    \end{tikzpicture}
}
    \caption{Architecture of the surrogate policy $\textcolor{mygreen}{\pi}$. The encoder maps the observation $O$ to the embedding, from which the decoder reconstructs the action distribution parameters.}
    \label{fig:surrogate_policy_architecture}
\end{figure}


% For any observation $O$, each $\SurrogatePolicy{i}$ should produce an action distribution that closely matches that of $\InitialPolicy$.
To train each surrogate policy $\SurrogatePolicy{i}$ to accurately replicate the action distribution, we define a loss function $\mathcal{L}$ optimized via gradient descent as introduced in Section~\ref{sec:learning}.
This loss, defined in Equation~\eqref{eq:car_loss}, is computed as the \gls{mse} between the parameters of the action distributions, obtained by applying the function $p$, which maps a distribution to its parameterization.
In practice, this function $p$ does not need to be explicitly defined, as neural policies directly output the parameters of the action distribution.

\begin{align}
    \min_{\Theta_i} \sum_{O \in \mathcal{O}}
    \mathrm{MSE}\Big(
        p \, \!\big(\InitialPolicy\left(O\right)\big),
        p \, \!\big(\textcolor{myred}{\pi'_{\Theta_i}}\left(O\right)\big)
    \Big)
    \label{eq:car_loss}
\end{align}

\subsection{Cluster extraction}

Once the surrogate policies are trained, their encoders are used to generate a point cloud from a set of observations.
This point cloud provides a sample of the \gls{ls} to be analyzed.
The proposed \gls{car} method adopts the deep clustering approach introduced in Section~\ref{sec:clustering}.
Under this approach, conceptually similar observations are expected to produce nearby \gls{lr}.
A \gls{cs} can therefore be viewed as a region of the \gls{ls} containing related \gls{lr}.
To identify these regions, a clustering algorithm is applied to the point cloud.

The clustering is performed independently for each surrogate policy.
For a surrogate policy $\SurrogatePolicy{i}$, applying its encoder $\textcolor{myblue}{f_{enc_{i}}}$ to a subset of observations $\mathbb{O} \subseteq \mathcal{O}$ yields the set of embeddings $\textcolor{myblue}{\mathcal{E}_{i}}$.
A clustering function $f_{\mathrm{clust}}$ then partitions the point cloud $\textcolor{myblue}{\mathcal{E}_{i}}$ into groups of nearby \gls{lr}, assigning each observation a cluster label and yielding the partition $\mathcal{C}_{i}$.
The same observation set $\mathbb{O}$ is used for every surrogate policy, so that the partitions $\mathcal{C}_{1}, \dots, \mathcal{C}_{n}$ are defined over the same observations and can be compared.
This comparison is the basis of the evaluation introduced in the following section.
Algorithm~\ref{alg:car} states the complete \gls{car} procedure.

\begin{algorithm}[t]
\caption{Clustering Agent Representation (CAR)}
\label{alg:car}

\Input{
\\
Initial policy $\InitialPolicy$\\
Observation dataset $\mathbb{O}$\\
Number of surrogate policies $n$\\
Surrogate policy architecture $\textcolor{myred}{\pi'}$ (encoder $\textcolor{myblue}{f_{enc}}$, decoder $f_{dec}$)\\
Learning rate $\eta$\\
Clustering function $f_{\mathrm{clust}}$
}

\Output{
\\
Trained surrogate policies $\textcolor{myred}{\pi'_{\Theta_1}}, \dots, \textcolor{myred}{\pi'_{\Theta_n}}$\\
Clusterings $\mathcal{C}_{1}, \dots, \mathcal{C}_{n}$
}

\AlgoRule

\tcp{Behavioral cloning dataset built from the target policy $\InitialPolicy$}
$\mathcal{D} \leftarrow \big\{ \big(O, p(\InitialPolicy(O))\big) \mid O \in \mathbb{O} \big\}$\nllabel{alg:car:dataset}\;

\BlankLine
\tcp{MSE loss between action-distribution parameters, as defined in Equation~\eqref{eq:car_loss}}
$\mathcal{L} \leftarrow$ \gls{mse}\nllabel{alg:car:loss}\;

\BlankLine
\tcp{Independently train $n$ surrogate policies by behavioral cloning}
\For{$i = 1, \dots, n$\nllabel{alg:car:train_loop}}{
    Initialize $\Theta_i$ randomly\;

    \BlankLine
    \tcp{Behavioral cloning of $\InitialPolicy$ via SGD, see Algorithm~\ref{alg:stochastic_gradient_descent}}
    $\Theta_i \leftarrow$ \FnSGD{$\mathcal{D}, \textcolor{myred}{\pi'}, \mathcal{L}, \eta, \Theta_i$}\nllabel{alg:car:sgd}\;
}

\BlankLine
\tcp{Cluster the latent representations of each surrogate policy}
\For{$i = 1, \dots, n$\nllabel{alg:car:cluster_loop}}{
    \tcp{Embeddings produced by the encoder of $\textcolor{myred}{\pi'_{\Theta_i}}$}
    $\textcolor{myblue}{\mathcal{E}_{i}} \leftarrow \big\{ \textcolor{myblue}{f_{enc_i}}(O) \mid O \in \mathbb{O} \big\}$\nllabel{alg:car:embed}\;

    \BlankLine
    \tcp{Partition of the embedded observations into clusters}
    $\mathcal{C}_{i} \leftarrow f_{\mathrm{clust}}(\textcolor{myblue}{\mathcal{E}_{i}})$\nllabel{alg:car:clust}\;
}

\BlankLine
\Return{\nllabel{alg:car:return}$\big(\textcolor{myred}{\pi'_{\Theta_1}}, \dots, \textcolor{myred}{\pi'_{\Theta_n}}\big), \big(\mathcal{C}_{1}, \dots, \mathcal{C}_{n}\big)$}

\end{algorithm}


Two choices govern this analysis without being part of the clustering algorithm itself.
The first is the choice of layer used to extract the representations.
Each hidden layer defines its own \gls{ls}, reflecting the information represented at that stage of the network.
\gls{car} can therefore be applied to the \gls{lr} extracted from any layer.
This allows the conceptual organization of the policy to be examined at different levels of the network.

The second is the number of clusters.
\gls{car} does not assume that a task admits a single correct number of concepts.
This choice instead determines the granularity at which the conceptual organization is examined.
A small number reveals broad, strategic concepts, while a larger number reveals finer, tactical ones.
The influence of both choices is studied in Chapter~\ref{chap:experiments}.


Figure~\ref{fig:space_invaders_projection} shows the outcome of this procedure for a surrogate policy trained on \SpaceInvadersName, using five clusters.
The clusters are computed in the original high-dimensional \gls{ls}.
The point cloud is a two-dimensional projection used only for visualization.
Each cluster gathers observations that correspond to a distinct, interpretable stage of a game:
\begin{itemize}
    \item \textbf{Cluster 1:} The start of a game, with few or no invaders neutralized.
    \item \textbf{Cluster 2:} The right-hand group of invaders eliminated.
    \item \textbf{Cluster 3:} The pursuit of the bonus saucer.
    \item \textbf{Cluster 4:} The end of a game, with a few fast invaders close to the ground.
    \item \textbf{Cluster 5:} The middle of a game, with the invaders descending but already thinned out.
\end{itemize}

\input{figures/clustering_agent_representation/space_invaders_projection.tex}

The position of an observation in the \gls{ls} thus reflects the stage of the environment it belongs to, suggesting a correspondence between the geometry of the \gls{ls} and the semantic properties of environment states.

This example illustrates what \gls{car} is expected to produce, but it relies on the visual inspection of a single surrogate policy, which is precisely the type of evidence whose limitations motivated this work.
Three aspects must therefore be evaluated before such an analysis can be trusted.
First, the surrogate policy must faithfully reproduce the initial policy so that its explanations remain relevant to the target policy.
Second, the recovered conceptual organization must be consistent across independently trained surrogate policies.
Finally, the clusters must capture genuine abstractions rather than simply reproducing structure already present in the observations or action distributions.

The following section introduces one evaluation metric for each of these aspects and details how it is computed.

